
\chapter{攻读硕士学位期间取得的研究成果}

\zihao{5}
\vspace{2pt}
一、已发表（包括已接受待发表）的论文，以及已投稿、或已成文打算投稿、或拟成文投稿的论文情况\underline{\textbf{（只填写与学位论文内容相关的部分）：}}

\begin{table}[h]
    \centering
    \zihao{-5}
    \renewcommand{\arraystretch}{2}
    \renewcommand{\tabularxcolumn}[1]{m{#1}}
    \begin{tabularx}{\textwidth}{|m{0.4cm}<{\centering}|m{1.8cm}<{\centering}|X<{\centering}|m{2.6cm}<{\centering}|m{1.4cm}<{\centering}|m{1.5cm}<{\centering}|m{1.2cm}<{\centering}|}
        \hline
        \textbf{序号} &
        \textbf{作者（全体作者，按顺序排列）} &
        \textbf{题\quad 目} &
        \textbf{发表或投稿刊物名称、级别} &
        \textbf{发表的卷期、年月、页码} &
        \textbf{与学位论文哪一部分（章、节）相关} &
        \textbf{被索引收录情况} \\
        \hline

        1 &
        \textbf{张三}，李四，王五 &
        An Example Paper Title Related to Chapter 3 &
        Example Conference / Journal Name &
        20XX年 &
        第三章 &
        示例 \\
        \hline

        & & & & & & \\[2.5em]
        \hline
        & & & & & & \\[2.5em]
        \hline
        & & & & & & \\[2.5em]
        \hline

    \end{tabularx}
\end{table}

\vspace{-1em}

注：在“发表的卷期、年月、页码”栏：

\vspace{-0.6em}
\indent 1 如果论文已发表，请填写发表的卷期、年月、页码；

\vspace{-0.6em}
\indent 2 如果论文已被接受，填写将要发表的卷期、年月；

\vspace{-0.6em}
\indent 3 以上都不是，请据实填写“已投稿”，“拟投稿”。

\vspace{-0.6em}
\indent 不够请另加页。

\vspace{0.6em}

二、与学位内容相关的其它成果（包括专利、著作、获奖项目等）

\vspace{0.5em}
\zihao{5}
\indent 专利：

\begin{tabular}[t]{@{}l@{}}
       1. 示例发明人甲、\textbf{张三}、示例发明人乙. 一种示例专利名称. 已受理，20XX.\\[0.3em]
       2. 示例发明人丙、\textbf{张三}. 一种示例专利名称及装置. 已公开，20XX.
    \end{tabular}
